\chapter{L11 -- Raggiungibilità e controllabilità all'origine}

\section{Idea della lezione}

In questa lezione introduciamo uno dei problemi centrali della teoria del controllo:

\begin{quote}
dato un sistema dinamico lineare, esiste un ingresso di controllo che permetta di portare lo stato dove vogliamo?
\end{quote}

Questa domanda può essere letta in due modi leggermente diversi.

\begin{itemize}
    \item \textbf{Raggiungibilità}: partendo dall'origine, quali stati posso raggiungere in un tempo finito?
    \item \textbf{Controllabilità all'origine}: partendo da uno stato iniziale assegnato, posso riportare il sistema all'origine in un tempo finito?
\end{itemize}

Nel caso dei sistemi lineari tempo-invarianti questi concetti si descrivono in modo molto chiaro attraverso:
\begin{itemize}
    \item la formula di evoluzione dello stato;
    \item uno spazio di stati raggiungibili;
    \item la matrice di raggiungibilità;
    \item nel continuo, anche il gramiano di raggiungibilità.
\end{itemize}

L’idea fondamentale da tenere a mente è questa:

\begin{quote}
la possibilità di muovere il sistema nello spazio di stato non dipende da un singolo ingresso scelto a mano, ma dalla struttura della coppia \((A,B)\).
\end{quote}

\section{Sistema lineare stazionario a tempo continuo}

Consideriamo un sistema lineare stazionario a tempo continuo
\begin{equation}
\Sigma_c:
\begin{cases}
\dot x(t)=Ax(t)+Bu(t),\\
y(t)=Cx(t)+Du(t),
\end{cases}
\qquad x(t_0)=x_0,
\label{eq:L11_sistema_tc}
\end{equation}
dove
\[
x(t)\in\mathbb{R}^n,\qquad
u(t)\in\mathbb{R}^m,\qquad
y(t)\in\mathbb{R}^p.
\]

La soluzione del sistema è
\begin{equation}
x(t)=e^{A(t-t_0)}x_0+\int_{t_0}^{t} e^{A(t-\tau)}Bu(\tau)\,d\tau.
\label{eq:L11_soluzione_tc}
\end{equation}

Questa formula separa in modo netto due contributi:
\begin{itemize}
    \item il termine
    \[
    e^{A(t-t_0)}x_0,
    \]
    che è l'\textbf{evoluzione libera} e dipende solo dallo stato iniziale;
    \item il termine
    \[
    \int_{t_0}^{t} e^{A(t-\tau)}Bu(\tau)\,d\tau,
    \]
    che è l'\textbf{evoluzione forzata} e dipende dall’ingresso.
\end{itemize}

\subsection{Riduzione al caso con stato iniziale nullo}

Valutiamo \eqref{eq:L11_soluzione_tc} all’istante finale \(t_f\):
\begin{equation}
x(t_f)=e^{A(t_f-t_0)}x(t_0)+\int_{t_0}^{t_f} e^{A(t_f-\tau)}Bu(\tau)\,d\tau.
\label{eq:L11_x_tf_tc}
\end{equation}

Portando a sinistra il termine di evoluzione libera otteniamo
\begin{equation}
x(t_f)-e^{A(t_f-t_0)}x(t_0)
=
\int_{t_0}^{t_f} e^{A(t_f-\tau)}Bu(\tau)\,d\tau.
\label{eq:L11_traslazione_tc}
\end{equation}

Questa formula è molto importante, perché ci dice che il problema

\begin{quote}
\emph{partendo da \(x(t_0)\), posso arrivare in \(x(t_f)\)?}
\end{quote}

è equivalente al problema

\begin{quote}
\emph{partendo dall’origine, posso arrivare nel vettore}
\[
\bar x(t_f):=x(t_f)-e^{A(t_f-t_0)}x(t_0)\, ?
\]
\end{quote}

Quindi, per studiare la raggiungibilità, basta concentrarsi sul caso \(x(t_0)=0\).

\section{Insieme degli stati raggiungibili dall'origine nel continuo}

Se \(x(t_0)=0\), la formula dello stato diventa
\begin{equation}
x(t_f)=\int_{t_0}^{t_f} e^{A(t_f-\tau)}Bu(\tau)\,d\tau.
\label{eq:L11_raggiungibile_tc}
\end{equation}

Questo porta alla seguente definizione.

\subsection*{Definizione}
Fissato \(t_f>t_0\), l’insieme degli stati raggiungibili dall’origine è
\begin{equation}
\mathcal{R}^{t_f}
=
\left\{
x\in\mathbb{R}^n \;:\;
\exists\,u(\cdot)\ \text{tale che}\
x=
\int_{t_0}^{t_f} e^{A(t_f-\tau)}Bu(\tau)\,d\tau
\right\}.
\label{eq:L11_def_R_tf}
\end{equation}

In altre parole, \(\mathcal{R}^{t_f}\) è l’insieme di tutti gli stati in cui posso portare il sistema all’istante \(t_f\), partendo dall’origine.

\subsection{Operatore di raggiungibilità}

Definiamo l’operatore lineare
\begin{equation}
\varphi:\mathcal{U}\to\mathbb{R}^n,
\qquad
\varphi(u)=\int_{t_0}^{t_f} e^{A(t_f-\tau)}Bu(\tau)\,d\tau,
\label{eq:L11_operatore_phi}
\end{equation}
dove \(\mathcal{U}\) è uno spazio di ingressi, ad esempio \(L^2([t_0,t_f],\mathbb{R}^m)\).

Allora
\begin{equation}
\mathcal{R}^{t_f}=\operatorname{Im}(\varphi).
\label{eq:L11_Rtf_immagine_phi}
\end{equation}

Quindi lo studio della raggiungibilità coincide con lo studio dell’immagine di un operatore lineare.

\begin{figure}[H]
\centering
\begin{tikzpicture}[>=Latex, node distance=4cm]
    \node[draw, circle, minimum size=1.7cm] (U) {$\mathcal U$};
    \node[draw, circle, minimum size=1.7cm, right of=U] (X) {$\mathbb{R}^n$};
    \draw[->, thick] (U) -- node[above] {$\varphi$} (X);
    \node[below=0.7cm of X] {\(\operatorname{Im}(\varphi)=\mathcal R^{t_f}\)};
\end{tikzpicture}
\caption{L’insieme raggiungibile dall’origine nel continuo è l’immagine dell’operatore di raggiungibilità.}
\label{fig:L11_phi}
\end{figure}

\section{Operatore aggiunto e gramiano di raggiungibilità}

Per analizzare meglio \(\operatorname{Im}(\varphi)\), introduciamo l’operatore aggiunto \(\varphi^\ast\).

\subsection{Prodotti scalari}

Sullo spazio di stato usiamo il prodotto scalare euclideo:
\[
\langle x_1,x_2\rangle_X=x_1^\top x_2.
\]

Sullo spazio degli ingressi \(\mathcal U=L^2([t_0,t_f],\mathbb{R}^m)\) usiamo
\[
\langle u_1,u_2\rangle_{\mathcal U}
=
\int_{t_0}^{t_f} u_1^\top(\tau)u_2(\tau)\,d\tau.
\]

Per definizione, \(\varphi^\ast:\mathbb{R}^n\to\mathcal U\) è l’operatore tale che
\[
\langle \varphi(u),x\rangle_X
=
\langle u,\varphi^\ast(x)\rangle_{\mathcal U}
\qquad
\forall u\in\mathcal U,\ \forall x\in\mathbb{R}^n.
\]

\subsection{Calcolo esplicito di \(\varphi^\ast\)}

Partiamo da
\[
\varphi(u)=\int_{t_0}^{t_f} e^{A(t_f-\tau)}Bu(\tau)\,d\tau.
\]
Allora
\[
\langle \varphi(u),x\rangle_X
=
x^\top\int_{t_0}^{t_f} e^{A(t_f-\tau)}Bu(\tau)\,d\tau
=
\int_{t_0}^{t_f} x^\top e^{A(t_f-\tau)}Bu(\tau)\,d\tau.
\]
Usando la proprietà di trasposizione,
\[
x^\top e^{A(t_f-\tau)}Bu(\tau)
=
u^\top(\tau)B^\top e^{A^\top(t_f-\tau)}x.
\]
Quindi
\[
\langle \varphi(u),x\rangle_X
=
\int_{t_0}^{t_f}
u^\top(\tau)\,B^\top e^{A^\top(t_f-\tau)}x\,d\tau
=
\langle u,\varphi^\ast(x)\rangle_{\mathcal U}.
\]
Si riconosce allora
\begin{equation}
\varphi^\ast(x)(\tau)=B^\top e^{A^\top(t_f-\tau)}x.
\label{eq:L11_phi_star}
\end{equation}

\subsection{Gramiano di raggiungibilità}

Consideriamo ora la composizione \(\varphi\varphi^\ast\):
\[
(\varphi\varphi^\ast)(x)
=
\int_{t_0}^{t_f}
e^{A(t_f-\tau)}B\Big(B^\top e^{A^\top(t_f-\tau)}x\Big)\,d\tau.
\]
Raccogliendo \(x\), otteniamo
\[
(\varphi\varphi^\ast)(x)
=
\left(
\int_{t_0}^{t_f}
e^{A(t_f-\tau)}BB^\top e^{A^\top(t_f-\tau)}\,d\tau
\right)x.
\]

Questo suggerisce di definire il \textbf{gramiano di raggiungibilità}
\begin{equation}
G_{\mathcal R}(t_f)
=
\int_{t_0}^{t_f}
e^{A(t_f-\tau)}BB^\top e^{A^\top(t_f-\tau)}\,d\tau.
\label{eq:L11_gramiano}
\end{equation}

Allora
\begin{equation}
\varphi\varphi^\ast(x)=G_{\mathcal R}(t_f)x.
\label{eq:L11_phi_phi_star}
\end{equation}

Poiché in dimensione finita vale
\[
\operatorname{Im}(\varphi)=\operatorname{Im}(\varphi\varphi^\ast),
\]
concludiamo che
\begin{equation}
\mathcal{R}^{t_f}
=
\operatorname{Im}(G_{\mathcal R}(t_f)).
\label{eq:L11_Rtf_gramiano}
\end{equation}

Quindi nel continuo lo spazio raggiungibile può essere studiato anche tramite il gramiano.

\section{Matrice di raggiungibilità nel continuo}

Introduciamo ora la classica matrice di raggiungibilità:
\begin{equation}
R=
\begin{bmatrix}
B & AB & A^2B & \cdots & A^{n-1}B
\end{bmatrix}.
\label{eq:L11_matrice_raggiungibilita}
\end{equation}

Il risultato fondamentale è che, per sistemi lineari stazionari,
\begin{equation}
\operatorname{Im}(G_{\mathcal R}(t_f))=\operatorname{Im}(R).
\label{eq:L11_gramiano_matrice}
\end{equation}

Quindi, nel continuo,
\begin{equation}
\mathcal{R}^{t_f}=\operatorname{Im}(R).
\label{eq:L11_Rtf_uguale_R}
\end{equation}

\subsection{Commento sul significato}

Questo è un fatto molto importante.  
Anche se il gramiano contiene un integrale e dipende apparentemente dall’istante finale \(t_f\), il sottospazio degli stati raggiungibili è lo stesso che si ottiene guardando la matrice finita
\[
R=
\begin{bmatrix}
B & AB & \cdots & A^{n-1}B
\end{bmatrix}.
\]

Quindi, in pratica, il problema della raggiungibilità si riduce a un semplice problema di rango.

\section{Sistema completamente raggiungibile}

\subsection*{Definizione}
La coppia \((A,B)\) si dice \textbf{completamente raggiungibile} se
\begin{equation}
\operatorname{rank}(R)=n.
\label{eq:L11_rank_full}
\end{equation}

In tal caso
\[
\operatorname{Im}(R)=\mathbb{R}^n,
\]
e quindi ogni stato è raggiungibile dall’origine in tempo finito.

\subsection{Interpretazione geometrica}

Se \(\operatorname{rank}(R)=r<n\), allora gli stati raggiungibili non riempiono tutto lo spazio di stato, ma solo un sottospazio di dimensione \(r\).

Per esempio:
\begin{itemize}
    \item se \(r=1\), gli stati raggiungibili stanno su una retta;
    \item se \(r=2\) in \(\mathbb{R}^3\), stanno su un piano;
    \item se \(r=n\), si riempie tutto \(\mathbb{R}^n\).
\end{itemize}

\begin{figure}[H]
\centering
\begin{tikzpicture}[>=Latex, scale=1.0]
    \draw[->] (0,0) -- (4,0) node[right] {$x_1$};
    \draw[->] (0,0) -- (0,3.5) node[above] {$x_2$};
    \draw[->] (0,0) -- (-1.6,-1.2) node[below left] {$x_3$};

    \filldraw[fill=gray!20, draw=black]
    (0.4,0.4) -- (2.8,0.7) -- (1.7,2.3) -- (-0.6,2.0) -- cycle;

    \node at (2.3,2.1) {\(\operatorname{Im}(R)\)};
\end{tikzpicture}
\caption{Se \(\operatorname{rank}(R)=2\) in \(\mathbb{R}^3\), gli stati raggiungibili formano un piano.}
\label{fig:L11_piano_raggiungibile}
\end{figure}

\section{Spazio raggiungibile da uno stato iniziale generico}

Se \(x_0\neq 0\), allora dallo sviluppo dello stato segue che
\[
x(t_f)=e^{A(t_f-t_0)}x_0+r,
\qquad r\in\mathcal{R}^{t_f}.
\]
Quindi l’insieme degli stati raggiungibili a partire da \(x_0\) è
\begin{equation}
\mathcal{R}^{t_f}_{x_0}=e^{A(t_f-t_0)}x_0+\mathcal{R}^{t_f}.
\label{eq:L11_R_x0}
\end{equation}

Questo insieme non è, in generale, un sottospazio vettoriale: è una \textbf{traslazione affine} dello spazio raggiungibile dall’origine.

\section{Caso discreto}

Passiamo ora al sistema lineare stazionario a tempo discreto:
\begin{equation}
\Sigma_d:
\begin{cases}
x(k+1)=Ax(k)+Bu(k),\\
y(k)=Cx(k)+Du(k).
\end{cases}
\label{eq:L11_sistema_td}
\end{equation}

Iterando la dinamica si ottiene
\begin{equation}
x(t)=A^t x_0+\sum_{k=0}^{t-1}A^{t-k-1}Bu(k).
\label{eq:L11_soluzione_td}
\end{equation}

Se \(x_0=0\), allora
\begin{equation}
x(t)=\sum_{k=0}^{t-1}A^{t-k-1}Bu(k).
\label{eq:L11_td_zero}
\end{equation}

Raccogliendo i termini rispetto agli ingressi:
\begin{equation}
x(t)=
\begin{bmatrix}
B & AB & A^2B & \cdots & A^{t-1}B
\end{bmatrix}
\begin{bmatrix}
u(t-1)\\
u(t-2)\\
\vdots\\
u(0)
\end{bmatrix}.
\label{eq:L11_td_fattorizzata}
\end{equation}

Definiamo allora la \textbf{matrice di raggiungibilità in \(t\) passi}
\begin{equation}
R_t=
\begin{bmatrix}
B & AB & A^2B & \cdots & A^{t-1}B
\end{bmatrix}.
\label{eq:L11_Rt}
\end{equation}

e il corrispondente insieme degli stati raggiungibili dall’origine in \(t\) passi:
\begin{equation}
\mathcal{R}_t=\operatorname{Im}(R_t).
\label{eq:L11_spazio_Rt}
\end{equation}

\subsection{Monotonia degli spazi raggiungibili}

Vale la catena crescente
\begin{equation}
\mathcal{R}_1\subseteq \mathcal{R}_2\subseteq \cdots \subseteq \mathcal{R}_t\subseteq \mathcal{R}_{t+1}\subseteq \cdots
\label{eq:L11_catena_Rt}
\end{equation}
perché tutto ciò che posso raggiungere in \(t\) passi posso ancora raggiungerlo in \(t+1\) passi, eventualmente scegliendo l’ultimo ingresso nullo.

Poiché lo spazio di stato ha dimensione \(n\), questa catena si stabilizza al più dopo \(n\) passi.

\subsection*{Definizione}
Il sistema discreto è \textbf{completamente raggiungibile} se esiste \(t\) tale che
\begin{equation}
\dim \mathcal{R}_t=n.
\label{eq:L11_completa_raggiungibilita_td}
\end{equation}

Per un sistema lineare tempo-invariante basta controllare la matrice
\[
R=
\begin{bmatrix}
B & AB & \cdots & A^{n-1}B
\end{bmatrix}.
\]

\section{Controllabilità all'origine}

Finora abbiamo chiesto: \emph{quali stati posso raggiungere partendo dall’origine?}  
Ora ribaltiamo il punto di vista: \emph{da quali stati iniziali posso tornare all’origine?}

\subsection{Definizione nel discreto}

Fissato \(t\), definiamo
\begin{equation}
\mathcal{C}_t=
\left\{
x_0\in\mathbb{R}^n \;:\;
\exists\,u(0),\dots,u(t-1)\ \text{tali che}\ x(t)=0
\right\}.
\label{eq:L11_Ct}
\end{equation}

Usando \eqref{eq:L11_soluzione_td}, la condizione \(x(t)=0\) diventa
\[
A^t x_0+\sum_{k=0}^{t-1}A^{t-k-1}Bu(k)=0.
\]
Quindi
\begin{equation}
-A^t x_0\in \mathcal{R}_t.
\label{eq:L11_zero_condition}
\end{equation}

\subsection{Raggiungibilità implica controllabilità all'origine}

Se il sistema è completamente raggiungibile in \(t\) passi, allora
\[
\mathcal{R}_t=\mathbb{R}^n.
\]
Quindi per ogni \(x_0\) vale automaticamente
\[
-A^t x_0\in\mathcal{R}_t,
\]
e dunque ogni stato iniziale può essere portato all’origine.

Perciò nel discreto vale
\begin{equation}
\text{raggiungibilità} \;\Longrightarrow\; \text{controllabilità all'origine}.
\label{eq:L11_ragg_impl_ctrl0_td}
\end{equation}

\subsection{Nel discreto il viceversa non vale in generale}

Nel discreto può accadere che
\[
\operatorname{Im}(A^t)\subseteq \mathcal{R}_t
\]
pur senza avere \(\mathcal{R}_t=\mathbb{R}^n\).

In questo caso:
\begin{itemize}
    \item ogni stato iniziale può essere portato all’origine;
    \item ma non ogni stato è raggiungibile partendo dall’origine.
\end{itemize}

Quindi nel discreto:
\begin{equation}
\boxed{
\text{raggiungibilità} \Longrightarrow \text{controllabilità all'origine},
\qquad
\text{ma non viceversa.}
}
\label{eq:L11_no_equiv_td}
\end{equation}

\section{Nel continuo raggiungibilità e controllabilità all'origine coincidono}

Nel continuo la situazione è diversa.  
Infatti la condizione per arrivare a zero all’istante \(t_f\) è
\[
0=e^{A(t_f-t_0)}x_0+\int_{t_0}^{t_f} e^{A(t_f-\tau)}Bu(\tau)\,d\tau,
\]
cioè
\[
-e^{A(t_f-t_0)}x_0\in \mathcal{R}^{t_f}.
\]
Ma la matrice esponenziale è sempre invertibile:
\[
e^{A(t_f-t_0)} \ \text{è invertibile}\qquad \forall t_f>t_0.
\]
Quindi il cambio di variabile
\[
z=e^{A(t_f-t_0)}x_0
\]
è biiettivo, e perciò

\begin{equation}
\boxed{
\text{nel continuo, controllabilità all'origine e raggiungibilità coincidono.}
}
\label{eq:L11_equiv_tc}
\end{equation}

\subsection{Significato intuitivo}

Nel continuo l’evoluzione libera è governata da una matrice invertibile.  
Per questo motivo, il problema
\begin{quote}
\emph{portare \(x_0\) a zero}
\end{quote}
è equivalente al problema
\begin{quote}
\emph{raggiungere un certo stato partendo dall’origine.}
\end{quote}

\section{Lo spazio raggiungibile è \(A\)-invariante}

Una proprietà strutturale molto importante è che lo spazio raggiungibile è \(A\)-invariante.

\subsection*{Proposizione}
Lo spazio
\[
\operatorname{Im}(R)
\]
è \(A\)-invariante, cioè
\[
x\in\operatorname{Im}(R)\quad\Longrightarrow\quad Ax\in\operatorname{Im}(R).
\]

\subsection*{Dimostrazione}
Se \(x\in\operatorname{Im}(R)\), allora esiste \(y\) tale che
\[
x=Ry=
\begin{bmatrix}
B & AB & A^2B & \cdots & A^{n-1}B
\end{bmatrix}y.
\]
Allora
\[
Ax=
\begin{bmatrix}
AB & A^2B & A^3B & \cdots & A^nB
\end{bmatrix}y.
\]
Per il teorema di Cayley--Hamilton, \(A^nB\) è combinazione lineare delle colonne
\[
B,\ AB,\ \dots,\ A^{n-1}B.
\]
Di conseguenza anche \(Ax\) è combinazione lineare delle colonne di \(R\), e quindi
\[
Ax\in\operatorname{Im}(R).
\]
\hfill\(\square\)

\subsection{Commento}

Questa proprietà sarà decisiva nella lezione successiva, quando riscriveremo il sistema separando:
\begin{itemize}
    \item la parte raggiungibile;
    \item la parte non raggiungibile.
\end{itemize}

Infatti il fatto che \(\operatorname{Im}(R)\) sia \(A\)-invariante permette di costruire una base di coordinate adatta a questa decomposizione.

\section{Esempio semplice di non completa raggiungibilità}

Consideriamo il sistema discreto
\[
A=
\begin{bmatrix}
1 & 0 & 0\\
0 & 0 & 0\\
1 & 0 & 0
\end{bmatrix},
\qquad
B=
\begin{bmatrix}
1\\
0\\
1
\end{bmatrix}.
\]
Allora
\[
AB=
\begin{bmatrix}
1\\
0\\
1
\end{bmatrix}=B,
\qquad
A^2B=AB=B.
\]
Quindi
\[
R=
\begin{bmatrix}
B & AB & A^2B
\end{bmatrix}
=
\begin{bmatrix}
1 & 1 & 1\\
0 & 0 & 0\\
1 & 1 & 1
\end{bmatrix}.
\]
Il rango è
\[
\operatorname{rank}(R)=1<3.
\]
Dunque il sistema non è completamente raggiungibile.

Tuttavia
\[
\operatorname{Im}(A)=\operatorname{span}\!\left\{\begin{bmatrix}1\\0\\1\end{bmatrix}\right\}
=
\operatorname{Im}(R_1).
\]
Questo significa che ogni \(-Ax_0\) appartiene già a \(R_1\), quindi il sistema è controllabile all’origine in un passo, pur non essendo completamente raggiungibile.

Questo esempio mostra in modo molto chiaro la differenza tra i due concetti nel discreto.

\section{Schema riassuntivo finale}

\begin{center}
\fbox{
\parbox{0.93\textwidth}{
\textbf{Tempo continuo}
\begin{itemize}
    \item
    \[
    x(t)=e^{A(t-t_0)}x_0+\int_{t_0}^{t}e^{A(t-\tau)}Bu(\tau)\,d\tau
    \]
    \item
    \[
    \mathcal{R}^{t_f}=\operatorname{Im}(R),
    \qquad
    R=\begin{bmatrix}B & AB & \cdots & A^{n-1}B\end{bmatrix}
    \]
    \item
    \[
    \operatorname{rank}(R)=n
    \iff
    \text{sistema completamente raggiungibile}
    \]
    \item
    \[
    \text{raggiungibilità}
    \iff
    \text{controllabilità all'origine}
    \]
\end{itemize}

\vspace{0.2cm}

\textbf{Tempo discreto}
\begin{itemize}
    \item
    \[
    x(t)=A^t x_0+\sum_{k=0}^{t-1}A^{t-k-1}Bu(k)
    \]
    \item
    \[
    \mathcal{R}_t=\operatorname{Im}(R_t),
    \qquad
    R_t=\begin{bmatrix}B & AB & \cdots & A^{t-1}B\end{bmatrix}
    \]
    \item
    \[
    \mathcal{R}_1\subseteq \mathcal{R}_2\subseteq\cdots
    \]
    \item
    \[
    \text{raggiungibilità}
    \Longrightarrow
    \text{controllabilità all'origine}
    \]
    ma in generale il viceversa \emph{non} vale.
\end{itemize}
}}
\end{center}

\section{Conclusione}

La lezione ha introdotto il primo vero strumento strutturale della teoria del controllo: lo \textbf{spazio raggiungibile}.  
Da qui in avanti l’idea sarà sempre la stessa:

\begin{quote}
capire se l’ingresso riesce a generare tutto lo spazio di stato, oppure solo una sua parte.
\end{quote}

Questa domanda verrà sviluppata ulteriormente nella prossima lezione, dove useremo l’\(A\)-invarianza dello spazio raggiungibile per costruire la \textbf{forma standard di raggiungibilità}.